%% =================================================================
%% Wei, Mei, Zhao, Xiao, Sun, Zhang, Guo.
%% "Nondestructively identifying the mechanical behavior of soft
%%  tissues using surface deformation with an explicit inverse
%%  approach."
%% Appl. Math. Modelling 134 (2024) 126--147.
%%
%% Reconstruction of page 16 (printed p. 141) as study notes.
%% Generated by: reproduce-multiple-pages-basic-tex (batch 13-16)
%%
%% Structure: Fig. 20 (ring results with 1% noise) + Table 3
%% (ring-structure error comparison) + Table 4 (regularization
%% factors for ring) + Fig. 21 (composite inclusion targets).
%% No equations, no prose.
%%
%% New structural variation: TWO tables on one page. Examples only
%% showed one table per page so far (Table 1 on page 11, Table 2
%% on page 14). This page tests whether the stable preamble and
%% the table template scale to multiple tables cleanly.
%% =================================================================

\documentclass[11pt]{article}

\usepackage[margin=1in]{geometry}
\usepackage{amsmath, amssymb}
\usepackage{mathrsfs}
\usepackage{bm}
\usepackage{xcolor}
\usepackage{fancyhdr}
\usepackage{booktabs}
\usepackage{multirow}
\usepackage{array}

\pagestyle{fancy}
\fancyhf{}
\lhead{\small Y.~Mei et al.}
\rhead{\small Applied Mathematical Modelling 134 (2024) 126--147}
\cfoot{\small 141 $\to$ \arabic{page}}
\renewcommand{\headrulewidth}{0.4pt}

\setcounter{page}{1}

\begin{document}

% ---------- Figure 20 placeholder ----------
\noindent
\begin{center}
\fbox{\parbox{0.92\textwidth}{\centering\vspace{1em}
\textbf{[Figure 20 --- placeholder]}\\[0.8em]
Six-panel comparison of \textbf{Explicit Inverse Method} (top row)
vs.\ \textbf{Implicit Inverse Method} (bottom row) on the
\textbf{Case 2 bilayer ring structure} with \textbf{1\,\% noise}
added, using 3, 6, and 9 surface displacement datasets.\\[0.4em]
\textbf{Explicit (top row)}: near-identical bilayer ring
reconstructions, visually indistinguishable from the Fig.~17
target. Colorbar ranges:
(a) 3 fields, SM $\in [0.129, 0.385]$\,MPa;
(b) 6 fields, SM $\in [0.132, 0.384]$\,MPa;
(c) 9 fields, SM $\in [0.130, 0.386]$\,MPa.
The explicit method tolerates 1\,\% noise essentially perfectly
on the ring problem, with all three panels landing within
$\pm 0.002$\,MPa of the 0.386 nominal outer layer.\\[0.4em]
\textbf{Implicit (bottom row)}: severe overshoot pattern.
Colorbar ranges:
(d) 3 fields, SM $\in [0.010, 1.251]$\,MPa --- \textbf{224\,\%
overshoot} above the 0.386 nominal outer layer; visually dominated
by green/yellow artifacts;
(e) 6 fields, SM $\in [0.010, 0.874]$\,MPa --- 126\,\% overshoot;
(f) 9 fields, SM $\in [0.010, 0.975]$\,MPa --- 152\,\% overshoot.
The implicit method's max is also non-monotone in this figure
($1.251 \to 0.874 \to 0.975$), continuing the pattern from Fig.~16
and Fig.~19.\\[0.3em]
\textbf{Takeaway}: as in the noise-free case (Fig.~19), the
implicit method is fundamentally failing on the ring problem. The
1\,\% noise addition does not change the qualitative picture ---
it is already broken without noise. See Table~3 for the
quantitative confirmation (implicit $\approx 41$\,\% vs.\ explicit
$\approx 5$\,\% average relative error).
\vspace{1em}}}
\end{center}
\vspace{0.3em}
\begin{center}
\textbf{Fig. 20.} The reconstructed results with varying 3, 6 and
9 surface displacement datasets (1\,\% noise is introduced,
Unit: MPa).
\end{center}

\vspace{1em}

% ---------- Table 3 ----------
\noindent\textbf{Table 3}\\[0.2em]
{\small Comparison of the average relative error in the shear
modulus distribution obtained by implicit inverse and explicit
inverse methods for the ring structure (1\,\% noise is introduced).}
\vspace{0.3em}

\begin{center}
\small
\begin{tabular}{@{}l ccc ccc@{}}
\toprule
Inverse method
  & \multicolumn{3}{c}{Explicit inverse method}
  & \multicolumn{3}{c}{Implicit inverse method} \\
\cmidrule(lr){2-4}\cmidrule(lr){5-7}
Number of surface displacement datasets
  & 3 & 6 & 9 & 3 & 6 & 9 \\
\midrule
Average relative error
  & 4.73\,\% & 5.38\,\% & 4.56\,\%
  & 42.41\,\% & 40.99\,\% & 41.23\,\% \\
\bottomrule
\end{tabular}
\end{center}

\vspace{1em}

% ---------- Table 4 ----------
\noindent\textbf{Table 4}\\[0.2em]
{\small The values of regularization factors used in Figs.~19 and 20.}
\vspace{0.3em}

\begin{center}
\small
\begin{tabular}{@{}ll ccc ccc@{}}
\toprule
\multicolumn{2}{l}{Inverse method}
  & \multicolumn{3}{c}{Explicit inverse method}
  & \multicolumn{3}{c}{Implicit inverse method} \\
\cmidrule(lr){3-5}\cmidrule(lr){6-8}
\multicolumn{2}{l}{Number of surface displacement datasets}
  & 3 & 6 & 9 & 3 & 6 & 9 \\
\midrule
\multirow{2}{*}{Ring structure}
  & no noise    & 0.00 & 0.00 & 0.00 & $1\times 10^{-12}$ & $1\times 10^{-12}$ & $1\times 10^{-12}$ \\
  & 1\,\% noise & 0.00 & 0.00 & 0.00 & $1\times 10^{-12}$ & $1\times 10^{-12}$ & $1\times 10^{-12}$ \\
\bottomrule
\end{tabular}
\end{center}

\vspace{1em}

% ---------- Figure 21 placeholder ----------
\noindent
\begin{center}
\fbox{\parbox{0.85\textwidth}{\centering\vspace{1em}
\textbf{[Figure 21 --- placeholder]}\\[0.8em]
\textbf{Target shear modulus distribution} for the \textbf{Case 3
composite with embedded inclusions}. Two sub-panels:\\[0.3em]
\textbf{(a) Sample 1}: a square blue background (soft tissue,
$\mu_\text{background} = 1.00$\,MPa) with a single red circular
inclusion at the center (tumor, $\mu_\text{inclusion} = 5.00$\,MPa).
Stiffness ratio 5:1, matching the Case~3 setup described on
page~140.\\[0.3em]
\textbf{(b) Sample 2}: same square background with \textbf{two}
red circular inclusions side-by-side in the upper half of the
domain.\\[0.3em]
Both panels share boundary conditions: the top edge has an array
of small downward-pointing arrows labeled ``displacement
measurement'' (these are the indentation-force application points
and where surface displacements are recorded for the inverse
problem); the bottom edge is marked with small circles/triangles
indicating fixed/vertical-direction constraints. Colorbar ticks
at $\{1.00, 2.00, 3.00, 4.00, 5.00\}$\,MPa, jet colormap.
\vspace{1em}}}
\end{center}
\vspace{0.3em}
\begin{center}
\textbf{Fig. 21.} Target shear modulus distribution of composite
with inclusions embedded (Unit: MPa).
\end{center}

% [page ends here — continues on page 17]

\end{document}
